\chapter{Background}
\label{cha:background}

This chapter provides the technical background for the two key challenges facing modern coding agents. Section~\ref{sec:bg_context} discusses the problem of providing accurate and complete context to LLMs operating on large code repositories, along with RepoGraph as a state-of-the-art context retrieval algorithm. Section~\ref{sec:bg_kvcache} addresses the latency bottleneck caused by redundant computation in agentic workflows, presenting CacheBlend as an approach to position-independent KV cache reuse. Finally, Section~\ref{sec:bg_motivation} motivates the need for a joint optimization that addresses both challenges in a unified system.

\section{Repository-Aware Context Retrieval}
\label{sec:bg_context}

% \myx{This section should not call "Introduction" (there can be only one intro in the whole paper). Maybe a better section title is "Existing Repository-Aware Context Retrieval Algorithms"?}

Modern code LLM agents have shown impressive capabilities, but most still struggle with incomplete or irrelevant context when working on large repositories. Many frameworks, from early tool-using agents to recent systems like SWE-Agent \cite{yang2024sweagent}, AutoCodeRover \cite{zhang2024autocoderover}, and Agentless \cite{xia2024agentlessdemystifyingllmbasedsoftware}, often feed the model only limited snippets of a codebase or generic summaries. Conventional agents often are able to interact with the entire codebase via searching keywords, files, and directories, in addition to some ``scrolling" mechanism for file viewing, {\colorGreen that is, having a constant-sized context window that can be moved up or down to have a view on an arbitrary code snippet of any file. SWE-agent's implementation, for example, consists of three commands: \verb|goto|, \verb|scroll_up|, and \verb|scroll_down|, allowing the agent to move the context window to any line of the target file \cite{yang2024sweagent}}. However, this does not prevent them from getting stuck in a local optimum \cite{ouyang2025repographenhancingaisoftware} \myx{what does "local optimum" refer to, better be specific}. {\colorGreen What this means is that the agent consistently has a limited view on the repository, resulting in possible false assumptions on code structure and implementation details.} This lack of essential context can lead to unstable performance and failures in code generation or bug-fixing (we refer to this as \textbf{Challenge 1}). For example, providing an agent with only a function implementation but not definitions of its referenced symbols (classes, functions, etc.) can omit crucial information and hurt the quality of patches\myx{A figure to illustrate this example would be good.} \mytodo{Add something like a referenced function is async but agent is unaware}. Conversely, naively supplying too much code (e.g., multiple flat files or irrelevant documentation/config) can overwhelm the model with irrelevant details, interfering with its reasoning \cite{zhang2024hierarchicalcontextpruningoptimizing, zhang2024inftybenchextendinglongcontext}. Previous studies have noted the challenge of balancing concise yet sufficient context given the finite context window of LLMs \cite{ouyang2025repographenhancingaisoftware}.

% Explorations have been made in an effort to mitigate Challenge 1. A common approach is retrieval-augmented generation (RAG), where the agent searches the repository for relevant snippets given an issue description. For instance, some systems treat the bug report as a query and use information retrieval (e.g.\ BM25 or dense embeddings) to pull in code segments likely related to the issue. Sliding window scans have also been used to ensure no part of a large file is missed by the context. Other frameworks rely on the LLM itself to navigate code: Agentless presents the entire repository structure (files and directories) to the model, letting it decide which files to open and read in a hierarchical fashion \cite{xia2024agentlessdemystifyingllmbasedsoftware}. In contrast, graph-based indexing methods pre-compute a knowledge graph of the repository (linking function definitions to references, etc.) and retrieve context via graph queries. These can provide stronger guidance, but maintaining a full repository graph (or vector index) is resource-intensive and struggles with continuously evolving codebases.

More broadly, code repositories present unique challenges for LLM-based agents due to their scale and interconnectivity. Real-world codebases often contain hundreds of thousands of lines of code spread across many files. This makes it infeasible to simply dump the entire repository into the LLM's context window. Instead, an agent must \emph{localize} the relevant portions of code for a given task. Traditional program analysis techniques, such as dependency graphs and call graphs, have long been used to understand large codebases, and recent LLM agents are beginning to incorporate such ideas. Early coding assistants relied on heuristic searches (e.g.\ grepping for error messages or function names) and would open files one by one in a Localize-then-Edit paradigm. More modern agents improve on this by indexing repository knowledge: for example, SWE-Search \cite{antoniades2025swesearchenhancingsoftwareagents} integrate a Monte Carlo Tree Search to explore the complex solution space, and CodePlan \cite{bairi2023codeplanrepositorylevelcodingusing} tracks incremental code changes with dependency analysis.

Despite these advancements, providing just the right context remains difficult. A key issue is that repository-level issue solving often requires following chains of references across multiple files. If any link in the chain is missing from the prompt, the model may fail to resolve the problem. Conversely, adding too much code can dilute the context. CoSIL \cite{jiang2025cosilsoftwareissuelocalization} \myx{Use this to support the sentence "Conversely, naively supplying too much code ..." in Sec. 2.1.} highlights this tension: it points out that existing methods struggle to give concise yet effective context coverage . CoSIL's solution is to dynamically construct a module-level call graph during the agent's search, expanding the graph by following function calls and pruning irrelevant branches.

RepoGraph \cite{ouyang2025repographenhancingaisoftware} takes a complementary approach by pre-building a complete repository graph before the agent starts working, offering a state-of-the-art solution to Challenge 1 by explicitly modeling the structure and dependencies of the repository. It constructs a comprehensive \emph{repository-level code graph} that captures relationships between code elements across different files. Each class, method, or function is a node, and the edges connect definitions to their usage or dependency. During a setup phase, RepoGraph parses all code files (using \texttt{tree-sitter}) to identify definitions (classes, functions, variables) and their references, filters out non-code files, and stores the resulting graph for fast querying. When the agent needs context for a particular function or error, it can query this graph to retrieve all related definitions and usage chains. In essence, RepoGraph provides an oracle that, given a symbol or filename, returns a subgraph containing all immediately relevant symbols. \mytodo{include a graphic example here}. By following these links, the agent is much less likely to overlook a relevant piece of code compared to linear search or limited file context. This approach has proven effective with GPT4/4o or Claude-3.5-Sonnet as backend: when plugged into Agentless or SWE-Agent, RepoGraph consistently boosted their success rates \cite{ouyang2025repographenhancingaisoftware}. This makes RepoGraph ideal as the foundation of our frontend context retrieval algorithm, whose specific usage will be discussed in closer detail in Chapter~\ref{cha:design}.

\mytodo{eval related insights, especially why comprehensive eval with diff APIs is important}

\section{Position-Independent KV Cache Reuse}
\label{sec:bg_kvcache}

% \myx{Same as the last section, you can rename the two subsections as "Existing solutions ..." and "Challenges ..." and exchange their ordering.}
% \myx{Also, chapter 3 might benefit from having a title similar to "Latency optimization through cached context"? As we want to transition from the latency challenge to the solution.}

Besides quality, the other major limitation of current code LLM agents is latency (\textbf{Challenge 2}). Many agentic frameworks send a large number of sequential requests to the LLM as it undergoes localization. This iterative process is significantly time-consuming, especially when the prompt in each request includes bulk context information as well as template tokens, and the model re-processes it from scratch every time. For instance, an agent might retrieve a large snippet of code from RepoGraph, feed it into the model to propose a fix, then in a subsequent step feed much of the same snippet again while testing or refining the fix. When no optimizations are prepared for this case, each of those calls recomputes attention over the repeated context tokens, undergoing the same KV cache prefill stage and wasting computation. This redundant computation is a key inefficiency that is shared by many agentic workflows.

At the core of this inefficiency is the key-value (KV) cache in transformer models, which stores intermediate attention states of processed tokens so that each new token need not recompute attention over the entire sequence. However, this cache normally resets between separate API calls. To reuse computation across calls, several optimized LLM serving engines have introduced caching and scheduling strategies. vLLM \cite{kwon2023vllm} improves throughput by batched execution and paged memory management of the KV cache, treating it as pageable memory. It can serve multiple generation requests in parallel and merge requests with identical prefixes so that the shared prefix is computed only once. However, even vLLM only reuses prefix computations when those prefixes exactly align from the start of the prompt. In an agent scenario, calls often share partial contexts that appear after some unique tokens, leading to low cache hit rates. SGLang \cite{zheng2024sglangefficientexecutionstructured} introduces RadixAttention, storing caches in a radix tree keyed by prefixes and scheduling requests to maximize cache hits, though it too relies on exact prefix matching. Another approach, Prompt Cache \cite{gim2024promptcachemodularattention}, allows explicitly marking reusable prompt segments and precomputing their KV states. While this yields substantial latency reductions on long-document QA tasks, it requires manual template setup and cannot adapt if templates change.

Cache optimization for LLM inference has recently taken a leap forward with CacheBlend \cite{yao2025cacheblendfastlargelanguage}. CacheBlend's key innovation is enabling \emph{position-independent} reuse of KV caches. It treats the prompt as composed of chunks, checks for each chunk if it has stored KV state, and when reassembling chunks selectively recomputes a small subset of tokens---specifically HKVD (High-KV-Deviation) tokens---to mitigate attention deviation and update positional encodings to reflect shifts in token positions, restoring the impact of cross-attention on the forward-attention matrix. The observation that HKVD tokens tend to remain deviated through different layers also makes it possible for efficient selection and KV recomputation of such tokens. This ``blending'' enables much higher cache hit rates on retrieval workflows, yielding 2--3$\times$ reduction in time-to-first-token and 3--5$\times$ throughput improvements, all without degrading output quality. CacheBlend was proved effective in multi-round QA scenarios through datasets such as 2WikiMQA and Musique \cite{ho2020constructingmultihopqadataset, trivedi2022musiquemultihopquestionssinglehop} \mytodo{maybe expand a bit on previous evaluations}. However, there remains a gap in understanding how CacheBlend performs in real-world agentic settings, where prompts can consist of long code context and complex repository structure information. Moreover, the original implementation assumes contiguous prefix alignment for cache lookup and allocates memory objects one at a time per cache chunk, which introduces scalability limitations under the high concurrency and variable prompt lengths characteristic of agentic workloads. In this work, we extend CacheBlend with non-contiguous cache hit support and enhanced memory allocation to address these limitations (see Chapter~\ref{cha:design}), and evaluate how the augmented system reduces latency under agentic use cases by integrating it into a RepoGraph-augmented agent scenario and measuring end-to-end latency against baseline systems like vLLM (with and without caching enabled).

\section{Motivation for Joint Optimization}
\label{sec:bg_motivation}

% \myx{Please ping me when you are done with this chapter.}

Having identified \textbf{Challenge 1} (poor patch quality due to incomplete context) and \textbf{Challenge 2} (slow generation due to long context) as two fundamental limitations of modern agentic systems, it is natural to consider designing an end-to-end system that aims at tackling both challenges, given the combined performance, latency, and cost benefits of a joint optimization. In the preceding sections, we surveyed state-of-the-art solutions for each respective challenge: RepoGraph \cite{ouyang2025repographenhancingaisoftware} for context retrieval and CacheBlend \cite{yao2025cacheblendfastlargelanguage} for inference latency. RepoGraph extracts repository-level context via pre-generating a structure tree containing symbol locations and a code graph containing call relations. CacheBlend enables position-independent KV cache full reuse by selectively recomputing a small subset of high-KV-deviation tokens to restore cross-attention. {\colorGreen In our integrated system, the aim is to combine the performance and latency advantages of each while mitigating the degradation of accuracy due to partially recomputed KV cache. The core concept that allows this is \textit{Cached Context}, which is a simple yet effective idea of chunking context blocks in prompts generated in the agentic workflow, such that the back-end KV cache reuse scheme (CacheBlend in our case) is able to exploit reuse opportunities more frequently. The reason why this is particularly effective in agentic settings is primarily because of the modular and repetitive nature of code agent context, which mostly consists of repository structure related info and function or class definition code blocks. Prompt templates, often containing format requirements and task-specific instructions, also will have tokens that are repeated upon every new request sent to the LLM. \textit{Cached Context} allows capturing these additional reuse opportunities to further speed up inference without loss of accuracy.}

 Mention LMCache's assumption on Long doc QA and not tailored for coding agent use